%============================================================================
% RESUMEN
%============================================================================
\begin{abstract}
Se estudió un enlace de telemetría industrial de baja velocidad afectado por
perturbaciones impulsivas de motores, contactores y variadores. El sistema se
simuló en banda base con BPSK, pulso rectangular,
$R_s=\SI{32}{\kilo\sym\per\second}$, ocho muestras por símbolo y
$f_s=\SI{256}{\kilo\samples\per\second}$, con ruido gaussiano de fondo, un tren
periódico de impulsos de amplitud controlada y un receptor con filtro acoplado.
Se compararon una referencia y tres niveles de amplitud mediante la señal
temporal, la PSD, la constelación, el diagrama de ojo y la BER. No hubo errores
mientras la amplitud efectiva del impulso, diluida por el filtro acoplado, fue
menor que la del símbolo; por encima, la BER alcanzó un piso de
$1{,}1\times10^{-2}$, independiente de la amplitud y reproducido por el modelo
analítico. La PSD no mostró cambios apreciables, por lo que las medidas
promediadas no bastan para diagnosticar ruido impulsivo. Además, una prueba
de laboratorio con 16QAM entre dos radios HackRF One mostró que la constelación
solo fue reconocible al aumentar la relación señal a ruido acercando las placas.
\end{abstract}

\begin{IEEEkeywords}
BPSK, ruido impulsivo, telemetría industrial, filtro acoplado, tasa de error de bit.
\end{IEEEkeywords}
